\documentclass[12pt,a4paper]{article}

\usepackage[margin=2.5cm]{geometry}
\usepackage{amsmath}
\usepackage{booktabs}
\usepackage{hyperref}
\usepackage{parskip}
\usepackage{microtype}
\usepackage{enumitem}
\usepackage{fancyhdr}
\setlength{\headheight}{14pt}
\newcommand{\degr}{\ensuremath{^\circ}}
\usepackage{titlesec}
\usepackage{xcolor}
\usepackage{listings}
\usepackage{caption}

% Header / footer
\pagestyle{fancy}
\fancyhf{}
\rhead{\small DSA-110 Continuum Imaging Pipeline}
\lhead{\small Technical Report}
\cfoot{\thepage}
\renewcommand{\headrulewidth}{0.4pt}

% Section title spacing
\titlespacing*{\section}{0pt}{1.4ex plus 0.5ex minus 0.2ex}{0.8ex plus 0.2ex}
\titlespacing*{\subsection}{0pt}{1.0ex plus 0.3ex minus 0.2ex}{0.5ex}

% Inline code style
\definecolor{codegray}{gray}{0.93}
\newcommand{\code}[1]{\colorbox{codegray}{\texttt{\small #1}}}

% -----------------------------------------------------------------------
\begin{document}

\begin{titlepage}
  \centering
  \vspace*{3cm}
  {\LARGE\bfseries The DSA-110 Continuum Imaging Pipeline\\[0.4em]
   Architecture, Data Flow, and Design Rationale\par}
  \vspace{1.5cm}
  {\large Technical Report\\[0.3em]
   Owens Valley Radio Observatory\par}
  \vspace{1.2cm}
  {\normalsize\itshape
   Compiled from source analysis of the\\
   \texttt{dsa110-contimg-legacy} repository\par}
  \vfill
  {\normalsize March 2026}
\end{titlepage}

% -----------------------------------------------------------------------
\section*{Abstract}
\addcontentsline{toc}{section}{Abstract}

The DSA-110 Continuum Imaging Pipeline is an automated, end-to-end software
system for producing calibrated radio continuum sky maps from the DSA-110
interferometric array at the Owens Valley Radio Observatory (OVRO).  The
pipeline ingests raw correlated visibilities stored in HDF5 format, applies
direction-independent calibration against reference catalogs, images and
mosaics the sky using the CASA toolkit, and performs aperture photometry to
identify transient and variable radio sources across multi-epoch observations.
This report describes the system architecture, the five sequential processing
stages, the three autonomous microservices that drive continuous operation, and
the key algorithmic decisions that reflect the instrument's specific operational
constraints.

% -----------------------------------------------------------------------
\section{Introduction}

The DSA-110 is a 110-dish interferometric array of 4.7\,m dishes operating at
L-band ($\sim$1.28--1.53\,GHz) at OVRO, California.  Its primary science
objective is the detection and sub-arcsecond localisation of fast radio bursts
(FRBs) through commensal wide-field survey observations
\cite{hallinan2019dsa}.  The continuum imaging pipeline operates in parallel
with the real-time transient detection backend, producing time-resolved sky
maps that support multi-epoch variability studies of persistent radio sources
and provide an independent verification channel for FRB host-galaxy
identification.

The pipeline is implemented in Python~3.10 and relies on CASA~6.6.3 (Common
Astronomy Software Applications) for the core interferometric operations ---
flagging, bandpass and gain calibration, deconvolution, and mosaicking.
Supporting astronomical libraries include \texttt{pyuvdata} (HDF5 data
ingest), \texttt{astroquery} (VizieR catalog access), \texttt{photutils}
(source detection), and \texttt{astropy} (coordinate handling, FITS I/O, WCS).
The system is structured around a modular, stage-based processing model with
three long-running asynchronous microservices that communicate via a
Redis-backed message queue, enabling fully autonomous continuous operation with
no human scheduling.

% -----------------------------------------------------------------------
\section{Pipeline Architecture}

The pipeline is organised into five sequential processing stages, orchestrated
by a central \code{PipelineOrchestrator} class.  A \code{ProcessingBlock}
abstraction encapsulates a single $\sim$5-minute observation chunk --- one
complete set of 16 HDF5 sub-band files --- as the fundamental unit of work.
Each stage is independently testable and can be skipped or replayed from
intermediate data products on disk.

\subsection{Stage 1 --- Data Ingestion}

Raw correlated visibilities are written by the DSA-110 correlator in UVH5
format (PyUVData's HDF5 dialect), with 16 sub-band files per 5-minute chunk
covering the full $\sim$250\,MHz bandwidth (\code{sb00} through \code{sb15}).
The \code{UnifiedMSCreationManager} reads each sub-band with
\code{run\_check=False} to bypass format validation, combines the 16 objects
in memory via the \code{UVData +=} operator, and writes a single CASA
Measurement Set (MS) spanning the full bandwidth.

A technically critical step is the correction of a systematic phase centre
error.  PyUVData's UVH5 reader silently ignores stored pointing coordinates
and defaults the phase centre to zenith (declination $= +90\degr$), which
would introduce a $\sim$37\degr\ pointing error in the output MS --- a
catastrophic but runtime-silent failure.  The fix reads the true phase centre
directly from the HDF5 header (\code{Header/phase\_center\_app\_ra},
\code{Header/phase\_center\_app\_dec}), overrides all PyUVData internal
attributes, and after writing explicitly patches the \code{REFERENCE\_DIR}
column of the MS \code{FIELD} subtable via \code{casatools.table}.

UVW coordinates are derived geometrically from a survey-grade ITRF antenna
position file (RevF survey) using
\code{set\_uvws\_from\_antenna\_positions()}, replacing an earlier erroneous
approach that scaled antenna positions to match raw UVW magnitudes.  An output
quality score (0--1) combines sub-band completeness (40\%), cross-file
consistency of timestamps and antenna counts (30\%), and integration-time
uniformity (30\%).

\subsection{Stage 2 --- Calibration}

Calibrator sources are selected automatically by querying the NVSS, FIRST,
TGSS, and VLASS catalogs via \texttt{astroquery}/VizieR for compact sources in
the flux range 3--15\,Jy at 1.4\,GHz within 1.5\degr\ of the survey
declination of $+16\degr$.  Results are cached locally to minimise repeated
network calls.  An NVSS cone search within the same radius additionally
produces a CASA component list used as a \code{startmodel} seed to accelerate
deconvolution.  Flux scaling between reference frequencies assumes a spectral
index of $\alpha = -0.7$.

CASA flagging (\code{rflag}, $3\sigma$ time and frequency cut-offs; \code{tfcrop};
shadow and autocorrelation clipping) is applied first.  Bandpass calibration
uses \code{solint=`inf'} (one solution per scan) with \code{solnorm=True};
gain calibration uses \code{solint=`1h'} in amplitude-and-phase mode
(\code{calmode=`ap'}).  Both tables are applied to the \code{CORRECTED\_DATA}
column of the MS via \code{applycal}.

\subsection{Stage 3 --- Imaging}

CASA \code{tclean} is run in multi-frequency synthesis (MFS) mode with the
\code{wproject} gridder, which corrects for the non-coplanar baseline effect
that arises from the DSA-110's $\sim$4\degr\ field of view at 1.4\,GHz --- a
regime in which the standard 2D FFT imaging approximation breaks down.  Key
parameters are summarised in Table~\ref{tab:tclean}.

\begin{table}[h!]
  \centering
  \caption{Default \texttt{tclean} parameters for DSA-110 imaging.}
  \label{tab:tclean}
  \begin{tabular}{ll}
    \toprule
    \textbf{Parameter} & \textbf{Value} \\
    \midrule
    \code{specmode}    & \code{mfs} \\
    \code{gridder}     & \code{wproject} (auto W-planes) \\
    \code{deconvolver} & \code{hogbom} \\
    \code{imsize}      & $4800 \times 4800$\,px \\
    \code{cell}        & $3''$/pixel \\
    \code{weighting}   & Briggs, robust $= 0.5$ \\
    \code{niter}       & 5000 \\
    \code{threshold}   & 1\,mJy \\
    \code{pbcor}       & True \\
    \bottomrule
  \end{tabular}
\end{table}

The grid of $4800 \times 4800$ pixels at $3''$/pixel covers $\approx4\degr
\times 4\degr$ and oversamples the $\sim$15$''$ synthesized beam by a factor
of five, trading storage for positional accuracy in multi-epoch source
matching.  Primary beam correction is applied and the result is exported to
FITS.  An image quality score is computed from the dynamic range, RMS noise,
and peak source flux.

\subsection{Stage 4 --- Mosaicking}

Approximately 12 individual 5-minute field images are combined per 60-minute
block using CASA \code{linearmosaic} with primary-beam weighting.  The output
mosaic spans $\approx 28800 \times 4800$ pixels, covering a $\sim$24\degr\
$\times$~4\degr\ strip along the survey declination, and is exported to FITS.

\subsection{Stage 5 --- Photometry}

Source detection and aperture photometry are performed on each mosaic FITS
image using \texttt{photutils}.  A $6''$ aperture ($\approx 2\times$ the
synthesized beam FWHM) with a $9''$--$15''$ sky annulus is used for all
sources.  Relative flux calibration is applied using an ensemble of up to 10
local NVSS reference sources within $15'$.  Photometric measurements --- flux
density, uncertainty, sky background, and epoch --- are written to a SQLite
(or optionally PostgreSQL) database keyed by source position.

% -----------------------------------------------------------------------
\section{Microservices and Continuous Operation}

For unattended autonomous operation, the pipeline stages are wrapped in three
asynchronous daemon services.

\textbf{HDF5 Watcher.}  A \code{watchdog.Observer} monitors the incoming HDF5
directory.  When a new file arrives, it is grouped with co-temporal sub-band
files into a \code{ProcessingBlock} and a \code{HDF5\_READY} event is
published to the Redis message queue.

\textbf{MS Processor.}  Consumes \code{HDF5\_READY} events and drives the
full five-stage pipeline for each block.  Processing results and error
conditions are published back to the queue.

\textbf{Variability Analyzer.}  Runs on a daily cron schedule (default
01:00~UTC).  Queries the photometry database over a 90-day lookback window
and flags as variability candidates any sources whose flux exhibits a
$\geq 5\sigma$ median absolute deviation (MAD) excursion on at least two
consecutive epochs.  Results are written to \code{variable\_candidates.csv}.

A Prometheus-compatible \texttt{aiohttp} dashboard exposes service health
metrics at \code{localhost:8080}.

% -----------------------------------------------------------------------
\section{Discussion}

Several design decisions reflect the specific operational constraints of the
DSA-110.  The choice of the \code{wproject} gridder over the computationally
cheaper \code{standard} gridder is required by the instrument's wide field of
view, which violates the coplanar-baseline approximation at L-band.  The
automated calibrator finder eliminates manual scheduling, which is essential
for a survey telescope that cycles through many pointings continuously without
human oversight.  The SQLite photometry store provides zero-infrastructure
persistence suitable for single-node operation, while the optional PostgreSQL
backend allows scaling to distributed deployments.

The most distinctive engineering challenge in the codebase is the phase centre
correction described in Section~2.1.  It illustrates a class of failure common
in scientific software pipelines: a library default that is silently incorrect
for the specific instrument's data format, producing no runtime exception but
yielding qualitatively wrong science output.  The fix required patching data at
three distinct levels simultaneously --- Python object attributes, PyUVData's
internal phase-centre catalog, and the on-disk CASA \code{FIELD} subtable ---
underscoring the importance of end-to-end integration tests that validate
output coordinate frames against known catalog positions rather than merely
confirming that processing steps complete without error.

% -----------------------------------------------------------------------
\section{Summary}

The DSA-110 Continuum Imaging Pipeline is a fully automated radio
interferometric reduction system spanning HDF5 ingest through multi-epoch
variability detection.  Its modular stage-based architecture, Redis-driven
microservice orchestration, and CASA-based calibration and imaging core make it
well suited to the continuous, unattended survey mode of the DSA-110.  Key
engineering contributions include a targeted correction of a systematic phase
centre error introduced by the upstream data-format library, a geometrically
consistent UVW recalculation from survey-grade antenna positions, and an
autonomous calibrator selection system that queries four archival radio
catalogs without human intervention.

% -----------------------------------------------------------------------
\begin{thebibliography}{9}

\bibitem{hallinan2019dsa}
  Hallinan, G.\ et al.\ 2019,
  \textit{BAAS}, 51, 261.
  The DSA-110: A Radio Survey Camera for FRB Localisation.

\bibitem{mcmullin2007casa}
  McMullin, J.~P.\ et al.\ 2007,
  \textit{ASP Conf.\ Ser.}, 376, 127.
  CASA Architecture and Applications.

\bibitem{offringa2014wsclean}
  Offringa, A.~R.\ et al.\ 2014,
  \textit{MNRAS}, 444, 606.
  WSClean: An Implementation of a Fast, Generic Wide-Field Imager for
  Radio Astronomy.

\bibitem{perley2017nvss}
  Condon, J.~J.\ et al.\ 1998,
  \textit{AJ}, 115, 1693.
  The NRAO VLA Sky Survey.

\end{thebibliography}

\end{document}
